% !TEX program = lualatex
% 14: デジタル表現2：画像・動画 / 圧縮 / デジタルデータの特徴

\documentclass[aspectratio=43,professionalfonts,handout,10pt]{beamer}
\usepackage{beamer_template}

\newcommand{\LectureNum}{14}
\newcommand{\LectureTitle}{デジタル表現2：画像 / 圧縮と品質}
\newcommand{\CourseName}{情報リテラシー}
\newcommand{\TextPages}{18-20}

\begin{document}
% --- Slide 1: title ---

\begin{frame}[shrink=5]{\CourseName\quad 第\LectureNum 回「\LectureTitle 」}
  {\small \TermName\quad \TeacherName\quad ／\quad 教科書 \TextPages}

  \vfill

  \begin{block}{今日の流れ}
    \begin{enumerate}
      \item 画像のデジタル化（Theory 18）
      \item 動画のデジタル表現（Theory 18）
      \item 圧縮とは何か（Theory 19）
      \item 可逆圧縮（Theory 19）
      \item 非可逆圧縮（Theory 19）
      \item デジタルデータの特徴（Theory 20）
    \end{enumerate}
  \end{block}
  \vfill

  \begin{block}{今日の目標}
    \begin{itemize}
      \item 画像・動画のデジタル表現，データ圧縮の方法，デジタルデータの特徴を説明できる。
    \end{itemize}
  \end{block}
  \vfill

  \begin{exampleblock}{キーワード}
    \small \textbf{画素（ピクセル）} ／ \textbf{解像度} ／ \textbf{階調} ／ \textbf{フレームレート（fps）} ／ \textbf{可逆圧縮} ／ \textbf{非可逆圧縮} ／ \textbf{圧縮率}
  \end{exampleblock}
\end{frame}

% --- Slide 2: 画像のデジタル化 (Theory 18 p.46) ---

\begin{frame}[shrink=15]{画像のデジタル化}
  \begin{block}{デジタル化の手順}
    \small
    \begin{itemize}
      \item ①標本化：画像を画素（ピクセル）という区画に分割し，各画素の値を取り出す
      \item ②量子化：取り出した値を段階値に変換（階調数 = $2^n$，$n$ = ビット数）
      \item ③符号化：量子化した数値を2進法で表す
      \item 解像度：ディスプレイ・印刷物の画素の密度。高いほど精細な画像になる
      \item 2値画像（1bit）・モノクロ256階調（8bit）・フルカラー（24bit = RGB各8bit）
    \end{itemize}
  \end{block}

  \vfill

  \begin{columns}[T]
    \begin{column}{0.48\textwidth}
      \begin{exampleblock}{色の表現}
        \small
        \begin{itemize}
          \item 加算混合（光）：R・G・B → 混ぜると白
          \item 減算混合（色材）：C・M・Y → 混ぜると黒に近い
        \end{itemize}
      \end{exampleblock}
    \end{column}
    \begin{column}{0.48\textwidth}
      \begin{block}{画像のデータ量}
        \small
        データ量（bit）\\
        = 横画素数 × 縦画素数\\
        \hspace{1em}× 1画素あたりのビット数
      \end{block}
    \end{column}
  \end{columns}

  \begin{alertblock}{演習}
    400×300画素，フルカラー（24bit）の画像のデータ量（バイト）を求めよ
  \end{alertblock}

  \begin{slidenote}
    \scriptsize 画素（ピクセル）はスマホ写真を最大拡大すると見える小さな正方形。フルカラー24bitはR・G・B各8ビット（256段階）×3色≈1,677万色（$2^{24}$）。加算混合（光）：R+G+B=白，減算混合（色材）：C+M+Y=黒に近い。【演習解答】$400\times300\times24\div8=360{,}000$ B。
  \end{slidenote}
\end{frame}

% --- Slide 3: 動画のデジタル表現 (Theory 18 p.47) ---

\begin{frame}[shrink=8]{動画のデジタル表現}
  \begin{block}{ポイント}
    \begin{itemize}
      \item \textbf{フレームレート（fps）}：動画で1秒間に表示するフレーム（コマ）数。多いほど滑らか
      \item 動画のデータ量 = 1枚あたりのデータ量 × フレームレート（fps）× 再生時間（秒）
    \end{itemize}
  \end{block}

  \vfill

  \begin{columns}[T]
    \begin{column}{0.48\textwidth}
      \begin{exampleblock}{インターレース方式}
        \small
        \begin{itemize}
          \item 画面を1行おきに走査して表示
          \item テレビ放送で使われてきた方式
        \end{itemize}
      \end{exampleblock}
    \end{column}
    \begin{column}{0.48\textwidth}
      \begin{block}{プログレッシブ方式}
        \small
        \begin{itemize}
          \item 画面を全行順番に走査して表示
          \item PCモニタ・現代ディスプレイの主流
        \end{itemize}
      \end{block}
    \end{column}
  \end{columns}

  \begin{alertblock}{演習}
    400×300画素，フルカラー（24bit）の画像を24fps，5秒の動画にした。データ量（B）を求めよ
  \end{alertblock}

  \begin{slidenote}
    \scriptsize インターレース：テレビ放送で使われてきた方式。プログレッシブ：PCモニタ・現代ディスプレイの主流。【演習解答】$400\times300\times24\div8=360{,}000$ B，$360{,}000\times24\times5=43{,}200{,}000$ B。
  \end{slidenote}
\end{frame}

% --- Slide 4: 圧縮とは何か (Theory 19 p.48) ---

\begin{frame}{圧縮とは何か}
  \begin{block}{圧縮の定義}
    \begin{itemize}
      \item 一定の手順を使ってデータを減らす技術を\textbf{圧縮}という
      \item \textbf{可逆圧縮}：圧縮されたデータを元と同じデータに戻せる
      \item \textbf{非可逆圧縮}：元と同じデータには戻せないが，人間が知覚しにくい情報を削除して圧縮率を上げる
    \end{itemize}
  \end{block}

  \vfill

  \begin{exampleblock}{COLUMN：ランレングス圧縮}
    \begin{itemize}
      \item 同じデータが続くとき，データと繰り返し回数で表す方法
      \item 例：「AAAABBBCC」→「4A3B2C」
      \item 256ピクセル（16×16）に4色のみ使用の場合，通常 $256\times8=2{,}048$ bit が必要なところ，大幅に削減できる
    \end{itemize}
  \end{exampleblock}

  \begin{slidenote}
    \scriptsize 圧縮の目的：保存容量の節約・通信時間の短縮。ランレングス圧縮は単純な繰り返しパターン（白黒画像・テキストなど）に有効。繰り返しが少ない画像には効果が薄い。
  \end{slidenote}
\end{frame}

% --- Slide 5: 可逆圧縮 (Theory 19 p.48) ---

\begin{frame}{可逆圧縮}
  \begin{block}{特徴}
    \begin{itemize}
      \item 圧縮されたデータを\textbf{元と同じデータに完全に戻せる}
      \item 文字データやプログラムなど，少しも欠けると正しく利用できないファイルに使用
    \end{itemize}
  \end{block}

  \vfill

  \begin{exampleblock}{主な形式}
    \begin{itemize}
      \item \textbf{ZIP・RAR}：複数ファイルをまとめて圧縮できる
      \item \textbf{PNG・GIF}：画像の可逆圧縮形式
    \end{itemize}
  \end{exampleblock}

  \begin{slidenote}
    \scriptsize ZIPはWindowsに標準搭載（右クリック→圧縮）。PNGはWebで広く使われる可逆圧縮画像形式（透過も対応）。GIFはアニメーション対応・256色制限あり。
  \end{slidenote}
\end{frame}

% --- Slide 6: 非可逆圧縮 (Theory 19 p.49) ---

\begin{frame}[shrink=20]{非可逆圧縮}
  \begin{block}{特徴}
    \small
    \begin{itemize}
      \item 元のデータには\textbf{完全に戻せない}
      \item 人間が知覚しにくい情報を削除して高い圧縮率を実現
      \item 圧縮率↑ → 品質↓ のトレードオフ
    \end{itemize}
  \end{block}

  \vfill

  \begin{exampleblock}{主な形式}
    \begin{itemize}
      \item \textbf{JPEG}（画像）：細かい色の差を間引く。写真向き
      \item \textbf{MP3}（音声）：人間が聞き取りにくい高音域などを削除
      \item \textbf{MPEG}（動画）：フレーム間の差分のみ記録し，変化のない部分を省略
    \end{itemize}
  \end{exampleblock}

  \begin{block}{圧縮率}
    \small 圧縮率 = 圧縮後のデータ量 ÷ 元のデータ量　（小さいほどよく圧縮されている）
  \end{block}

  \begin{alertblock}{演習}
    \small
    \begin{enumerate}
      \item[(1)] 400×300画素，フルカラー（24bit），30fps，10秒の動画の無圧縮データ量（MB）を求めよ（1MB=1,048,576B）
      \item[(2)] (1)の動画をMPEG圧縮（圧縮率0.05）した。圧縮後のデータ量（MB）を求めよ
      \item[(3)] (1)の動画を10Mbpsの回線で10秒間にリアルタイム送信したい。必要な最大圧縮率を求めよ
    \end{enumerate}
  \end{alertblock}

  \begin{slidenote}
    \scriptsize JPEG：スマホ写真の標準形式（高圧縮）。保存を繰り返すたびに劣化する。MP3：CDの音楽データを約1/10に圧縮できる。MPEG：YouTubeやDVDで使われる。【演習解答】(1) $400\times300\times24\div8=360{,}000$ B，$360{,}000\times30\times10\div1{,}048{,}576\approx102.9$ MB。(2) $102.9\times0.05\approx5.1$ MB。(3) $10\times10^6\times10\div8=12{,}500{,}000$ B，$12{,}500{,}000\div108{,}000{,}000\approx0.1157$。
  \end{slidenote}
\end{frame}

% --- Slide 7: デジタルデータの特徴 (Theory 20 p.50) ---

\begin{frame}[shrink=15]{デジタルデータの特徴}
  \begin{block}{プラス面}
    \small
    \begin{itemize}
      \item \textbf{多種多様な形態の情報を統合できる}：文字・音・画像をすべて0と1で統合して記録できる
      \item \textbf{情報が劣化しない}：複製・伝送を繰り返しても品質が変わらない
      \item \textbf{情報の加工が容易}：コピー・編集・加工をソフトウェアで簡単に行える
      \item \textbf{管理しやすい}：大量のデータを小さなメディアに保存・検索・整理できる
    \end{itemize}
  \end{block}

  \begin{exampleblock}{マイナス面}
    \small
    \begin{itemize}
      \item \textbf{情報が失われる}：量子化誤差（A/D変換時に微妙な情報が失われる）
      \item \textbf{著作権侵害}：劣化なくコピーできるため無断複製・拡散しやすい
      \item \textbf{情報流出}：小さなメディアに大量データを保存できるため流出しやすい
      \item \textbf{データが膨大}：高品質化に伴いデータ量が増大し続ける
    \end{itemize}
  \end{exampleblock}

  \begin{slidenote}
    \scriptsize アナログテープは複製するたびに劣化するが，デジタルは何度コピーしても品質が変わらない。この特性は便利な反面，違法コピーのしやすさにもつながる（著作権問題）。USBメモリ1本に数万件の個人情報が入るため，紛失・盗難のリスクも高い。
  \end{slidenote}
\end{frame}

% --- チェックリスト ---

\begin{frame}{まとめ・チェックリスト}
  \begin{block}{自己チェック（できたら $\checkmark$ をつけよう）}
    \begin{itemize}
      \item[$\square$] 画像のデジタル化（画素・解像度・階調・RGB）とデータ量計算を理解した
      \item[$\square$] 動画のフレームレートとデータ量計算を理解した
      \item[$\square$] 可逆圧縮と非可逆圧縮の違いと圧縮率の意味を理解した
      \item[$\square$] デジタルデータのプラス面・マイナス面を理解した
    \end{itemize}
  \end{block}

  \vfill

  {\small 質問があれば授業後またはTeamsで受け付けます。}
\end{frame}

\end{document}
